% CESC 410L -- Lab 0. Figures only, per the handout.
%
% DELIBERATE EXCEPTION: this report carries its own preamble rather than
% sharing the one the homework solutions use. In order of weight:
%   1. A handed-in .tex has to compile wherever the reader puts it. A shared
%      preamble is reached by a relative path that resolves only at this
%      file's exact depth, so the submitted copy would not build.
%   2. The packaging check rejects a repository path on a non-comment line,
%      and an \input of the course macros is exactly that.
%   3. A report is prose plus figures under a title block. It wants no running
%      head, no answer box and no problem header -- which is most of what the
%      shared preamble exists to provide.
% The overlap is four package lines and a margin. That is the whole price.

\documentclass[11pt]{article}

\usepackage[margin=1in]{geometry}
\usepackage{graphicx}
\usepackage{caption}
\usepackage{float}

% Look for figures beside this file first, then one level up. The second entry
% lets a copy of this document sitting in a subfolder still find figs/.
\graphicspath{{./}{../}}

\setlength{\parindent}{0pt}

\title{CESC 410L --- Lab 0\\[0.3em]
       \large Getting Started with Python Programming for DSP}
\author{Gatlin Nelson}
\date{\today}

\begin{document}
\maketitle

Base frequency $\mathrm{Freq\_1} = 4$\,Hz, from seed \texttt{myID = 6127}.

In every figure the upper panel shows the sparse samples (dots) over the dense
waveform (line); the lower panel shows the magnitude spectrum, annotated with
the complex DFT value at each peak.

\begin{figure}[H]
  \centering
  \includegraphics[width=0.72\textwidth]{figs/lab0_sinusoids_fig1.png}
  \caption{Single sinusoid. $A = 1$, $f = 4$\,Hz, $\phi = \pi/4$.}
\end{figure}

\begin{figure}[H]
  \centering
  \includegraphics[width=0.72\textwidth]{figs/lab0_sinusoids_fig2.png}
  \caption{Single sinusoid. $A = 0.5$, $f = 8$\,Hz, $\phi = \pi/3$.}
\end{figure}

\begin{figure}[H]
  \centering
  \includegraphics[width=0.72\textwidth]{figs/lab0_sinusoids_fig3.png}
  \caption{Superposition of the two sinusoids above, summed as signals.}
\end{figure}

\begin{figure}[H]
  \centering
  \includegraphics[width=0.72\textwidth]{figs/lab0_sinusoids_fig4.png}
  \caption{Superposition of three sinusoids built from parameter vectors:
           $A = [1,\ 0.5,\ 0.2]$, $f = [4,\ 8,\ 12]$\,Hz,
           $\phi = [\pi/4,\ \pi/3,\ \pi/9]$.}
\end{figure}

\end{document}
